\section{Verified-intervention dataset}
\label{sec:dataset}

\subsection{Collection protocol}

The dataset is constructed from two sources. The randomized online study
(R2-C1) generates controller decisions and outcomes under six policies
across at least 32 task instances and 120 fault episodes. The offline
replay dataset uses R1 development traces to construct eligible decision
points with action availability, cost, subsequent evidence, and outcome.
Where the logging policy provides adequate propensity support, inverse-
propensity or doubly robust estimates are used for off-policy evaluation;
unsupported counterfactual actions remain explicitly unknown.

\subsection{Transition record}

Each verified-intervention transition records:

\begin{itemize}
    \item the observable history and calibrated belief $b_t$ before the action;
    \item all actions available under policy and resource constraints, with
          eligibility masks;
    \item the selected action $a_t$ and its selection propensity
          $p(a_t \mid b_t)$;
    \item the action cost vector $\mathbf{c}(a_t)$, the resulting state
          difference, and the immediate postcondition;
    \item the final scientific outcome under independent adjudication;
    \item replayed counterfactual outcomes for alternative actions where
          replay validity permits;
    \item ambiguity flags, verifier disagreements, and missing evidence;
    \item the source split, deduplication cluster identifier, and
          eligibility for each training view.
\end{itemize}

Sampling intentionally oversamples rare fault--action combinations for
training diversity but retains prevalence weights for population-level
evaluation. Training balance is not mistaken for deployment prevalence.

\subsection{Split architecture}

The intervention dataset uses the same six split regimes as the R1
benchmark: development (pilot-only), unseen task templates, unseen injector
implementations, unseen fault families, unseen fault compositions, and an
external transfer arm. The primary R2 evaluation uses unseen fault
compositions and clean tasks. Additional evaluation sets include verifier
substitution (a different verifier implementation evaluating the same
trajectories) and environment/provider shift (tasks executed under a
changed runtime or model version).

\subsection{Dataset gates}

The verified-intervention dataset is considered frozen when: (1)~each
primary recovery action has adequate independent episode coverage across at
least two task families; (2)~no single task template contributes more than
a predeclared share of any action class; (3)~at least two adjudicator
implementations or a human audit cover high-reward cases; (4)~duplicate and
near-duplicate clusters are isolated by split; and (5)~ambiguous and failed
transitions remain in the dataset for audit. Failure to meet a gate results
in a narrower dataset claim rather than silently relaxing the criterion.
